\documentclass[a4paper,12pt]{article}
\usepackage[utf8]{inputenc}
\usepackage[T1]{fontenc}
\usepackage[ngerman]{babel}
\usepackage{amsmath}
\usepackage{amssymb}
\usepackage{enumitem}
\usepackage{array}
\usepackage{geometry}
\geometry{a4paper, margin=2.5cm}
\usepackage{fancyhdr}
\usepackage{parskip}

% Punkteangabe rechtsbündig am Ende eines Aufgabenteils
\newcommand{\pkt}[1]{\hfill\textit{(#1\,P)}}
% Kopf einer Aufgabe: \aufgabe{Nummer}{Titel}{Gesamtpunkte}
\newcommand{\aufgabe}[3]{%
  \bigskip\par\noindent\textbf{Aufgabe #1: #2}\hfill\textbf{#3 Punkte}%
  \par\nobreak\medskip}

\pagestyle{fancy}
\fancyhf{}
\lhead{\small 61211 Analysis -- Probeklausur}
\rhead{\small Matr.-Nr.: \makebox[3cm]{\dotfill}}
\cfoot{\small Seite \thepage}
\renewcommand{\headrulewidth}{0.4pt}

\begin{document}
\thispagestyle{empty}

\begin{center}
  {\LARGE\bfseries Probeklausur}\\[0.4em]
  {\large 61211 Analysis}\\[0.6em]
  {\normalsize FernUniversität in Hagen $\cdot$ Fakultät für Mathematik und Informatik}
\end{center}

\vspace{1.2em}

\noindent
\begin{tabular}{@{}ll@{}}
  Name: & \makebox[7cm]{\dotfill}\\[1.0em]
  Vorname: & \makebox[7cm]{\dotfill}\\[1.0em]
  Matrikelnummer: & \makebox[7cm]{\dotfill}\\
\end{tabular}

\vspace{1.5em}

\begin{center}
  \renewcommand{\arraystretch}{1.5}
  \begin{tabular}{|l|c|c|c|c|c|c|c||c|}
    \hline
    \textbf{Aufgabe} & 1 & 2 & 3 & 4 & 5 & 6 & 7 & \textbf{$\Sigma$}\\
    \hline
    max.\ Punkte & 20 & 12 & 12 & 12 & 16 & 14 & 14 & 100\\
    \hline
    erreicht & & & & & & & & \\
    \hline
  \end{tabular}
\end{center}

\vspace{1.0em}

\noindent\textbf{Bearbeitungshinweise}
\begin{itemize}[leftmargin=1.5em, itemsep=2pt]
  \item Die Bearbeitungszeit beträgt \textbf{120 Minuten}. Es sind alle 7 Aufgaben zu bearbeiten.
  \item Zugelassenes Hilfsmittel: ein handbeschriebenes DIN-A4-Blatt (auch beidseitig) mit eigenen Notizen.
  \item Verwenden Sie keinen Bleistift und keinen Rotstift.
  \item Formulieren Sie Ihre Lösungen in vollständigen Sätzen und begründen Sie Ihre Schritte. Verweise auf Sätze des Kurstextes dürfen benutzt werden.
\end{itemize}

\newpage

% =====================================================================
\aufgabe{1}{Multiple Choice}{20}

Kreuzen Sie an, ob die jeweilige Aussage \emph{wahr} oder \emph{falsch} ist. Für jede
richtig beantwortete Aussage erhalten Sie $2$ Punkte, für jede unbeantwortete Aussage
$1$ Punkt und für jede falsch beantwortete Aussage $0$ Punkte.

\medskip
\renewcommand{\arraystretch}{1.5}
\noindent\begin{tabular}{@{}c@{\hspace{0.6em}}p{11.4cm}@{\hspace{0.8em}}c@{\hspace{1.0em}}c@{}}
  & & \small Wahr & \small Falsch\\[0.2em]
  (a) & Für alle $x\in\mathbb{R}^n$ gilt $\|x\|_2 \le \|x\|_1$. & $\square$ & $\square$\\
  (b) & Jede Cauchyfolge in $(\mathbb{R}^n,\|\cdot\|_2)$ konvergiert. & $\square$ & $\square$\\
  (c) & Auf $\mathbb{R}^n$ sind die Normen $\|\cdot\|_1$ und $\|\cdot\|_\infty$ \emph{nicht} äquivalent. & $\square$ & $\square$\\
  (d) & Das stetige Bild einer zusammenhängenden Menge ist zusammenhängend. & $\square$ & $\square$\\
  (e) & Jede auf ganz $\mathbb{R}$ stetige Funktion $f:\mathbb{R}\to\mathbb{R}$ nimmt ihr Maximum an. & $\square$ & $\square$\\
  (f) & Ist $f:\mathbb{R}^2\to\mathbb{R}$ in $a$ partiell differenzierbar, so ist $f$ in $a$ total differenzierbar. & $\square$ & $\square$\\
  (g) & Ist $f:\mathbb{R}^n\to\mathbb{R}$ in $a$ differenzierbar mit $\operatorname{grad} f(a)\neq 0$, so zeigt $\operatorname{grad} f(a)$ in die Richtung des stärksten Anstiegs von $f$ in $a$. & $\square$ & $\square$\\
  (h) & Ist die Hesse-Matrix einer zweimal stetig differenzierbaren Funktion $f$ in einem kritischen Punkt $a$ indefinit, so besitzt $f$ in $a$ kein lokales Extremum. & $\square$ & $\square$\\
  (i) & Konvergiert eine Folge stetiger Funktionen $f_k:[a,b]\to\mathbb{R}$ punktweise gegen $f$, so gilt stets $\displaystyle\int_a^b f_k(t)\,dt \longrightarrow \int_a^b f(t)\,dt$. & $\square$ & $\square$\\
  (j) & Besitzt ein stetiges Vektorfeld $f:G\to\mathbb{R}^n$ auf einem Gebiet $G$ eine Stammfunktion, so sind seine Kurvenintegrale wegunabhängig. & $\square$ & $\square$\\
\end{tabular}

% =====================================================================
\aufgabe{2}{Gradient und Richtungsableitung}{12}

Sei $f:\mathbb{R}^3\to\mathbb{R}$, $f(x,y,z):=x^2y+yz^2-xz$, und $a:=(1,-1,2)^{T}$.
\begin{enumerate}[label=\alph*)]
  \item Begründen Sie kurz, dass $f$ auf ganz $\mathbb{R}^3$ differenzierbar ist, und berechnen Sie $\operatorname{grad} f(a)$. \pkt{4}
  \item Berechnen Sie die Richtungsableitung $D_v f(a)$ für $v:=\tfrac13(2,-2,1)^{T}$. Weisen Sie zuvor nach, dass $v$ ein zulässiger Richtungsvektor ist. \pkt{4}
  \item In welche Richtung wächst $f$ im Punkt $a$ am stärksten, und wie groß ist dort der maximale Wert einer Richtungsableitung $D_w f(a)$ (über alle $\|w\|_2=1$)? \pkt{4}
\end{enumerate}

% =====================================================================
\aufgabe{3}{Cauchy-Schwarz als Werkzeug}{12}

Sei $n\in\mathbb{N}$ und $c\in\mathbb{R}$ fest. Betrachten Sie alle Vektoren
$x=(x_1,\dots,x_n)\in\mathbb{R}^n$ mit $\sum_{k=1}^n x_k = c$.
\begin{enumerate}[label=\alph*)]
  \item Beweisen Sie mit der Cauchy-Schwarzschen Ungleichung (Satz~1.1.8), dass für jedes solche $x$ gilt
  \[ \sum_{k=1}^n x_k^2 \ge \frac{c^2}{n}. \]
  \pkt{5}
  \item Bestimmen Sie alle $x$, für die in a) Gleichheit gilt. \emph{Hinweis:} Überlegen Sie, wann in der Cauchy-Schwarzschen Ungleichung Gleichheit eintritt. \pkt{4}
  \item Werten Sie das Ergebnis für $n=5$, $c=10$ konkret aus: Wie lautet der kleinstmögliche Wert von $\sum x_k^2$, und für welches $x$ wird er angenommen? \pkt{3}
\end{enumerate}

% =====================================================================
\aufgabe{4}{Stammfunktion und Wegunabhängigkeit}{12}

Gegeben sei das Vektorfeld
\[ f:\mathbb{R}^2\longrightarrow\mathbb{R}^2,\qquad
   f(x,y):=\begin{pmatrix} 2xy+y\cos(xy)\\ x^{2}+x\cos(xy)\end{pmatrix}. \]
\begin{enumerate}[label=\alph*)]
  \item Prüfen Sie, ob $f$ ein Gradientenfeld ist, und bestimmen Sie gegebenenfalls eine Stammfunktion $F$. \pkt{6}
  \item Berechnen Sie $\displaystyle\int_W f\cdot d\!\begin{pmatrix}x\\y\end{pmatrix}$ für die Kurve $W:=S(a,c)+S(c,b)$ mit
  $a=\binom{0}{0}$, $c=\binom{1}{0}$, $b=\binom{1}{\pi/2}$. Dabei bezeichne $S(p,q)$ die geradlinige Verbindungsstrecke von $p$ nach $q$. \pkt{4}
  \item Warum liefert jede andere stückweise glatte Kurve von $a$ nach $b$ denselben Wert? \pkt{2}
\end{enumerate}

% =====================================================================
\aufgabe{5}{Satz über implizite Funktionen}{16}

Betrachtet werde $F=\,^{t}(F_1,F_2):\mathbb{R}^4\to\mathbb{R}^2$ mit
\[ F_1(x,y,u,v):=u^2+v^2+x-y-1,\qquad F_2(x,y,u,v):=u+v+xy-2, \]
und der Punkt $c=\,^{t}(1,1,1,0)$.
\begin{enumerate}[label=\alph*)]
  \item Zeigen Sie $F(c)=0$ und dass sich die Gleichung $F(x,y,u,v)=0$ in einer Umgebung von $c$ eindeutig \glqq nach $(u,v)$ auflösen\grqq{} lässt, d.h.\ es gibt eine stetig differenzierbare Funktion $g=\,^{t}(g_1,g_2)$ auf einer Umgebung $U$ von $(1,1)$ mit $(u,v)=g(x,y)$ und $g(1,1)=(1,0)$. \pkt{7}
  \item Berechnen Sie die Jacobi-Matrix $g'(1,1)$. \pkt{6}
  \item Lässt sich $F(x,y,u,v)=0$ nahe $c$ auch eindeutig \glqq nach $(x,y)$ auflösen\grqq? Begründen Sie kurz (ohne die auflösende Funktion anzugeben). \pkt{3}
\end{enumerate}

% =====================================================================
\aufgabe{6}{Zusammenhang und ganzzahlige Werte}{14}

Sei $(X,\|\cdot\|)$ ein normierter Raum, $\emptyset\neq M\subseteq X$ eine
\emph{zusammenhängende} Teilmenge und $f:M\to\mathbb{R}$ stetig.
\begin{enumerate}[label=\alph*)]
  \item Beweisen Sie: Nimmt $f$ nur ganzzahlige Werte an (d.h.\ $f(M)\subseteq\mathbb{Z}$), so ist $f$ konstant. \pkt{9}
  \item Beweisen oder widerlegen Sie: Die Aussage aus a) bleibt richtig, wenn man die Voraussetzung \glqq $M$ zusammenhängend\grqq{} streicht. \pkt{5}
\end{enumerate}
Benutzen Sie ausschließlich Sätze aus Kapitel~2.

% =====================================================================
\aufgabe{7}{Fourierreihe}{14}

Sei $f:\mathbb{R}\to\mathbb{R}$ die $2\pi$-periodische Fortsetzung von
\[ f(x)=\cos\!\Big(\frac{x}{2}\Big),\qquad x\in[-\pi,\pi]. \]
\begin{enumerate}[label=\alph*)]
  \item Zeigen Sie, dass $f$ gerade ist, und berechnen Sie die Fourierkoeffizienten. Weisen Sie nach, dass
  \[ a_k=\frac{4(-1)^{k+1}}{\pi\,(4k^{2}-1)}\ (k\in\mathbb{N}_0),\qquad b_k=0\ (k\in\mathbb{N}). \]
  \pkt{8}
  \item Begründen Sie mit dem Darstellungssatz~5.4.9, dass $f$ in jedem Punkt durch seine Fourierreihe dargestellt wird, und leiten Sie durch Auswertung an geeigneten Stellen die Reihensummen
  \[ \sum_{k=1}^{\infty}\frac{1}{4k^{2}-1}=\frac{1}{2}\qquad\text{und}\qquad
     \sum_{k=1}^{\infty}\frac{(-1)^{k+1}}{4k^{2}-1}=\frac{\pi}{4}-\frac{1}{2} \]
  her. \pkt{6}
\end{enumerate}

\vfill
\begin{center}\small Viel Erfolg!\end{center}

\end{document}
